% ============================================================================
\section{Sentiment Classification}
\label{sec:sentiment}
% Source: assignments/02_sentiment/index.html
% ============================================================================

% ----------------------------------------------------------------------------
\subsection{Building intuition}

Mini-reviews: $(-1)$ \emph{not good}; $(-1)$ \emph{pretty bad}; $(+1)$ \emph{good plot}; $(+1)$ \emph{pretty scenery}.
$\phi(x)$ = word counts and
$\text{Loss}_{\text{hinge}}(x,y,\bvec w)=\max\{0,\,1-\bvec w\cdot\phi(x)\,y\}$.

\problempart{1a}{Four SGD steps}
\question{Run one SGD update per review, in the order given, with $\eta=0.1$ from $\bvec w=\bvec 0$
(gradient $=0$ when the margin is exactly 1). Give the final weights of
(\emph{pretty, good, bad, plot, not, scenery}).}
\expect{The weight vector in that order, e.g.\ $[0.1,0.2,0.3,0.4,0.5,0.6]$.}
\begin{solution}
\todo{show each update, then the final vector}
\end{solution}

\problempart{1b}{A non-separable dataset}
\question{Dataset: $(-1)$ \emph{bad}; $(+1)$ \emph{good}; $(+1)$ \emph{not bad}; $(-1)$ \emph{not good}.
Prove that no linear classifier $f_{\bvec w}(x)=\sign(\bvec w\cdot\phi(x))$ on word-count features achieves
zero error, then propose one extra feature that fixes this.}
\expect{A short proof ($\sim$3--5 sentences) and a viable additional feature.}
\begin{solution}
\todo{answer}
\end{solution}

% ----------------------------------------------------------------------------
\subsection{Predicting movie ratings}

Predictor $\sigma(\bvec w\cdot\phi(x))$ with $\sigma(z)=(1+e^{-z})^{-1}$, rating $y\in[0,1]$.
(No Wolfram Alpha or similar.)

\problempart{2a}{Squared loss}
\question{Write $\text{Loss}(x,y,\bvec w)$ for squared loss on a single data point.}
\expect{A mathematical expression (may use $\sigma$).}
\begin{solution}
\todo{answer}
\end{solution}

\problempart{2b}{Its gradient}
\question{Compute $\nabla_{\bvec w}\text{Loss}(x,y,\bvec w)$ in terms of $p=\sigma(\bvec w\cdot\phi(x))$.}
\expect{A mathematical expression for the gradient.}
\begin{solution}
\todo{answer}
\end{solution}

\problempart{2c}{Vanishing gradients}
\question{For one data point with $y=1$, under what conditions on $\bvec w$ does $\norm{\nabla_{\bvec w}\text{Loss}}$
become arbitrarily small? Can it be exactly zero ($\norm{\bvec w}$ large but finite)?}
\expect{1--2 sentences on the conditions; 1--2 sentences on whether it can be exactly zero.}
\begin{solution}
\todo{answer}
\end{solution}

% ----------------------------------------------------------------------------
\subsection{Sentiment classification}
\codepart{2.3a--2.3d}
\texttt{extract\_word\_features}, \texttt{learn\_predictor} (SGD on hinge loss; train error $<4\%$, validation
error $<30\%$), \texttt{generate\_dataset}, \texttt{extract\_character\_features}.

\problempart{3e}{Character $n$-grams}
\question{Run the linear predictor with \texttt{extract\_character\_features} for several $n$. Which $n$
gives the smallest validation error, and why is it nearly as good as word features? Construct a one-sentence
review for which character $n$-grams probably beat word features, and explain why.}
\expect{A paragraph ($\sim$4--6 sentences) with the best $n$ and the reason; a one-sentence review with an explanation.}
\begin{solution}
\todo{table of validation error vs.\ $n$ (from notebook \S2.3e) and discussion}
\end{solution}

% ----------------------------------------------------------------------------
\subsection{Toxicity classification and maximum group loss}

$\phi(x)=[1,d,t]$: bias, demographic mention $d\in\{0,1\}$, toxic words $t\in\{0,1\}$; label $y\in\{-1,1\}$
($1$ = toxic). Groups are defined by $(y,d)$. Zero-one loss;
$\TrainLossMax(\bvec w)=\max_g\TrainLossG(\bvec w)$.

\begin{center}
\begin{tabular}{rrrr}
\toprule
$y$ & $d$ & $t$ & \# points\\
\midrule
$-1$&0&0&63\\ $-1$&0&1&27\\ $-1$&1&0&7\\ $-1$&1&1&3\\
$1$&0&0&3\\ $1$&0&1&7\\ $1$&1&0&27\\ $1$&1&1&63\\
\midrule
\multicolumn{3}{r}{Total} & 200\\
\bottomrule
\end{tabular}
\qquad
Classifier D: $\bvec w=[-0.1,1,0]$\quad Classifier T: $\bvec w=[-0.1,0,1]$
\end{center}

\problempart{4a}{Behaviour of D and T}
\question{Describe in words how each classifier behaves.}
\expect{For each classifier, an ``if and only if'' statement for when $f_{\bvec w}(x)=1$.}
\begin{solution}
\todo{answer}
\end{solution}

\problempart{4b}{Losses of classifier D}
\question{Compute D's average loss, its per-group losses, and its maximum group loss.}
\expect{The average loss, the completed table below, and the maximum group loss.}
\begin{solution}
\begin{center}
\begin{tabular}{lcc}
\toprule
 & $y=1$ & $y=-1$\\
\midrule
$d=1$ & \todo{} & \todo{}\\
$d=0$ & \todo{} & \todo{}\\
\bottomrule
\end{tabular}
\end{center}
Average loss: \todo{}. Maximum group loss: \todo{}.
\end{solution}

\problempart{4c}{Losses of classifier T; comparison}
\question{Compute the same three quantities for T. Which classifier has the lower average loss, and which
has the lower maximum group loss?}
\expect{The average loss, completed table, and maximum group loss, plus both comparisons.}
\begin{solution}
\begin{center}
\begin{tabular}{lcc}
\toprule
 & $y=1$ & $y=-1$\\
\midrule
$d=1$ & \todo{} & \todo{}\\
$d=0$ & \todo{} & \todo{}\\
\bottomrule
\end{tabular}
\end{center}
Average loss: \todo{}. Maximum group loss: \todo{}. Comparison: \todo{}.
\end{solution}

\problempart{4d}{Fairness and the mixed objective}
\question{$\TrainLossNew(\bvec w)=\lambda\,\TrainLossAvg(\bvec w)+(1-\lambda)\,\TrainLossMax(\bvec w)$, with
$\lambda\in[0,1]$. Use the ethical frameworks from the \emph{Algorithms and Distribution} module.}
\begin{enumerate}[label=(\roman*)]
  \item $\lambda=1$: describe the optimal classifier and the framework it aligns with. \expect{1--2 sentences, justified by the framework.}
  \begin{solution}\todo{answer}\end{solution}
  \item $\lambda=0$: same question. \expect{1--2 sentences, justified by the framework.}
  \begin{solution}\todo{answer}\end{solution}
  \item $\lambda=0$ with a new training set in which group $(y=1,d=1)$ has size 1 (the other groups have
        65, 65, 69). How would you factor group size into $\TrainLossMax$? \expect{The adjustment, justified in 1--2 sentences.}
  \begin{solution}\todo{answer}\end{solution}
  \item What $\lambda$ would you deploy on a real platform that cares equally about average and maximum group
        loss? Discuss the trade-offs with reference to the frameworks. \expect{3--5 sentences.}
  \begin{solution}\todo{answer}\end{solution}
\end{enumerate}

% ----------------------------------------------------------------------------
\subsection{K-means clustering}

$\phi(x_1)=[0,0]$, $\phi(x_2)=[4,0]$, $\phi(x_3)=[6,0]$, $\phi(x_4)=[11,0]$.

\problempart{5a}{2-means by hand}
\question{Run 2-means to convergence from (i) $\mu_1=[0,0],\ \mu_2=[11,0]$ and (ii) $\mu_1=[0,0],\ \mu_2=[4,0]$.
Give the assignments $z$, centres $\mu$ and final total loss.}
\expect{For each initialisation: centres and assignments at every step, and the final loss.}
\begin{solution}
\todo{answer for both initialisations}
\end{solution}

\problempart{5b}{Implementing k-means}
\codepart{2.5b}

\problempart{5c}{Scaling invariance}
\question{If every dimension of the data and the initial centroids is scaled by the same non-zero factor, are
the same clusters guaranteed? What if only some dimensions are scaled?}
\expect{For each sub-question: a short explanation if the answer is yes, or a concrete numerical counterexample if no.}
\begin{solution}
\todo{answer}
\end{solution}
